AI alignment is often framed as the problem of bringing a model's objectives, outputs, preferences, or behavior into accord with human values and intentions. In long-horizon human-AI systems that maintain memory, connect multiple models and tools, and modify external state, however, a major accident pathway arises when a locally valid output is promoted into display, execution, persistence, publication, or reinjection after losing the provenance, scope, time, uncertainty, representational conditions, or authority that justified it. The human supervising that promotion is also not an ideal adjudicator with unlimited attention, memory, time, or independence.

This paper connects two design perspectives, the \textbf{safe promotion of bounded artifacts} and the \textbf{finite-bandwidth human operator}, and expands the implementation unit of AI safety from the model alone to the sociotechnical system that generates, inspects, promotes, executes, verifies, and returns artifacts to memory. Each artifact retains not only content but also provenance, scope, time, generation or measurement conditions, missingness and uncertainty, permitted uses, and authority state. Capability is separated from authority. Authority is not granted wholesale to a system or session, but assigned to an atomic target use as a one-shot, expiring, and revocable authorization. Between approval and execution, state, evidence, budget, permission, and stop conditions are revalidated at use time.

The paper assumes that corrigibility, AI control, action alignment, least privilege, attribute-based access control, zero trust, and runtime governance already address authority, stopping, and external enforcement. Its novelty claim is not authority control itself. Rather, it treats epistemic artifacts, external execution, independent verification, memory, publication, and reinjection within a single promotion graph, and joins false promotion, false suppression, and finite-bandwidth human control in one audit vocabulary.

The paper advances five principal recommendations. First, evaluate AI safety as a property of systems rather than models alone. Second, institutionally separate capability evaluation from the granting of authority. Third, treat human-in-the-loop not as the presence of a human but as a conformance state satisfying referential clarity, local sufficiency, intent ownership, evidence calibration, legitimate authority, decision capacity, temporal freshness, and downstream understanding. Fourth, measure not only false promotion but also false suppression of legitimate candidates, pre-observation loss, and missing audit paths. Fifth, make the responsibilities of model developers, agent-platform developers, deploying organizations, evaluators, and standards and policy bodies explicit at each authority-promotion boundary.

Boundary-preserving alignment does not assume a perfectly benevolent model. It begins from the premise that models, evaluators, humans, and organizations can all fail, and requires a structure in which an error is not immediately promoted into an irreversible outcome. It is not a solution to the alignment problem. It is an implementable governance layer for preventing unresolved alignment problems from being coupled directly to unbounded authority.

\textbf{Keywords:} AI safety, AI alignment, agentic systems, bounded artifacts, bounded authority, human oversight, authority promotion, least privilege, provenance, finite-bandwidth operator, auditability, recoverability
